\subsection{Desempeño cinemático de la torreta (cañón) del prototipo}

\subsubsection{Descripción de la prueba}

Esta prueba valida la capacidad de movimiento angular de la torreta en sus dos ejes de control (\textit{pitch} y \textit{yaw}), de acuerdo con los requerimientos de seguimiento y apuntamiento del sistema.

\paragraph{Objetivos mínimos}
\begin{itemize}
    \item Alcanzar una velocidad angular en \textit{pitch} de $15\ \text{deg/s}$ sobre una apertura de $30^\circ$.
    \item Alcanzar una velocidad angular en \textit{yaw} de $40\ \text{deg/s}$ sobre un recorrido de $180^\circ$.
    \item Lograr en \textit{pitch} un ángulo máximo de elevación de $50^\circ$ respecto de la horizontal.
\end{itemize}

\subsubsection{Criterios de aprobación}

\begin{table}[H]
\centering
\small
\begin{tabular}{p{8.8cm} p{4.1cm}}
\toprule
\textbf{Métrica} & \textbf{Criterio de aprobación} \\
\midrule
Velocidad angular media en \textit{pitch} & $\geq 15\ \text{deg/s}$ \\
Apertura útil en \textit{pitch} & $\geq 30^\circ$ \\
Velocidad angular media en \textit{yaw} & $\geq 40\ \text{deg/s}$ \\
Recorrido útil en \textit{yaw} & $\geq 180^\circ$ \\
Ángulo máximo de elevación en \textit{pitch} & $\geq 50^\circ$ \\
Error de posicionamiento angular final (ambos ejes) & $\leq 2^\circ$ \\
\bottomrule
\end{tabular}
\caption{Criterios de aprobación de la prueba cinemática de torreta.}
\label{tab:criterios_prueba_canon}
\end{table}

\subsubsection{Procedimiento de ejecución}

\begin{enumerate}
    \item Fijar el prototipo en banco de ensayo para desacoplar efectos de tracción.
    \item Calibrar cero mecánico y horizontal de referencia para ambos ejes.
    \item Ejecutar barridos de \textit{pitch} sobre al menos $30^\circ$ y medir tiempo total de maniobra.
    \item Ejecutar barridos de \textit{yaw} sobre al menos $180^\circ$ y medir tiempo total de maniobra.
    \item Repetir cada barrido tres veces para estimar repetibilidad.
    \item Medir ángulo máximo de elevación alcanzable en \textit{pitch}.
    \item Calcular velocidad angular media y error de posicionamiento al final de cada maniobra.
\end{enumerate}

\subsubsection{Resultados}

\begin{table}[H]
\centering
\small
\begin{tabular}{p{8.8cm} p{2.2cm} p{2.8cm}}
\toprule
\textbf{Métrica} & \textbf{Resultado} & \textbf{Cumple} \\
\midrule
Velocidad angular media en \textit{pitch} [deg/s] & --- & --- \\
Apertura útil en \textit{pitch} [deg] & --- & --- \\
Velocidad angular media en \textit{yaw} [deg/s] & --- & --- \\
Recorrido útil en \textit{yaw} [deg] & --- & --- \\
Ángulo máximo de elevación en \textit{pitch} [deg] & --- & --- \\
Error angular final [deg] & --- & --- \\
\bottomrule
\end{tabular}
\caption{Resultados medidos de la prueba cinemática de torreta (a completar).}
\label{tab:resultados_prueba_canon}
\end{table}
